\subsection{子群上的诱导结构}

引理 4. 设 G 是有限维李群，\(\Omega\) 是 G 中 \(e\) 的对称开邻域，H 是 \(\Omega\) 的一个包含 \(e\) 的子集，并满足：条件 \(x \in \mathrm{H}\)、\(y \in \mathrm{H}\)、\(xy^{-1} \in \Omega\) 推出 \(xy^{-1} \in \mathrm{H}\)。设 \(r \in \mathrm{N}_{\mathbb{K}}\)。对所有 \(x \in \mathrm{H}\)，令 \(\mathfrak{h}_x\) 为所有满足下述性质的 \(a \in \mathrm{T}_x(\mathrm{G})\) 的集合：存在 K 中 0 的一个开邻域 I 及一个映射 \(f\)，它是从 I 到 G 的 \(\mathrm{C}^r\) 类映射，使得 \(f(0) = x, f(I) \subset H\)，且 \((\mathrm{T}_0f)(1) = a\)。

\begin{enumerate}
\def\labelenumi{(\roman{enumi})}
\item
  令 \(\mathfrak{h}_e = \mathfrak{h}\)。则 \(\mathfrak{h}\) 是 \(\mathrm{L}(\mathrm{G})\) 的李子代数，并且在 \(\mathrm{Ad}_{\mathrm{L}(\mathrm{G})}(\mathrm{H})\) 下不变。
\item
  \(\mathfrak{h}_x = x\mathfrak{h} = \mathfrak{h}x\) 对所有 \(x\in \mathrm{H}\) 成立（其中 \(x\mathfrak{h}\) 和 \(\mathfrak{h}x\) 在 \(\mathrm{T(G)}\) 中计算）。
\item
  设 \(\mathbf{V}\) 是 \(\mathbf{C}^r\) 类流形，\(v_0\) 是 \(\mathbf{V}\) 中的一点，且有一个映射 \(f\)，它是 \(\mathbf{C}^r\) 类的、从 \(\mathbf{V}\) 到 \(\mathbf{G}\) 的映射，满足 \(f(v_0) = e\) 且 \(f(\mathbf{V}) \subset \mathbf{H}\)。对于每个李子群芽 \(\mathbf{H}'\)，它属于 \(\mathbf{G}\)，且满足 \(\mathfrak{h}, f(v) \in \mathbf{H}'\)；当 \(v\) 充分接近 \(v_0\) 时这一关系成立。
\item
  对于每个李子群芽 \(\mathbf{H}'\)，它是 \(\mathbf{G}\) 的且李代数为 \(\mathfrak{h}\)，\(\mathbf{H}' \cap \mathbf{H}\) 是 \(e\) 在 \(H'\) 中的一个邻域。
\item
  对所有 \(x \in \mathbf{H}\) 及所有 \(a \in \mathfrak{b}_x\)，存在 K 中 0 的一个开邻域 I 及一个从 I 到 G 的映射 \(f\)，它属于 \(\mathbf{C}^{\omega}\) 类，使得 \(f(0) = x, f(I) \subset H, (T_0f)(1) = a\)。显然 \(\mathrm{K}\mathfrak{b} = \mathfrak{b}\)，并且 \(x\mathfrak{b}_y z = \mathfrak{b}_{xyz}\) 在 \(x, y, xy, xyz\) 属于 H 时成立。这推出 (ii)，并推出 \(\mathfrak{b}\) 在 \(\mathrm{Ad}_{\mathrm{L(G)}}(\mathrm{H})\) 下不变。
\end{enumerate}

设 \(a_1, a_2\) 属于 \(\mathfrak{h}\)。设 I 是 K 中 0 的一个开邻域，\(f_1, f_2\) 是从 I 到 G 的 \(\mathrm{C}^r\) 类映射，满足 \(f_j(0) = e, f_j(I) \subset H\)，\((\mathrm{T}_0 f_j)(1) = a_j (j = 1, 2)\)。定义 \(f: I \to G\) 为 \(f(\lambda) = f_1(\lambda)f_2(\lambda)\)。则 \(f\) 属于 \(\mathrm{C}^r\) 类且 \(f(0) = e\)。必要时缩小 I，可使 \(f(I) \subset H\)。另一方面，从 \(\mathrm{T}_e(G) \times \mathrm{T}_e(G)\) 到 \(\mathrm{T}_e(G)\) 的映射 \((g, g') \to gg'\) 在单位元处的切映射是加法；因此 \((\mathrm{T}_0 f) = a_1 + a_2\)。所以 \(a_1 + a_2 \in \mathfrak{h}\)，且 \(\mathfrak{h}\) 是 L(G) 的向量子空间。由于 \(x\mathfrak{h}x^{-1} = \mathfrak{h}\) 对所有 \(x \in H\) 都成立，\((\mathrm{Ad}f_1(\lambda)).a_2 \in \mathfrak{h}\) 对所有 \(\lambda \in I\) 都成立。根据 § 3 第12号命题44，映射 \(\lambda \mapsto \mathrm{Ad}f_1(\lambda)\) 在 0 处的切映射是映射 \(\lambda \mapsto \mathrm{ad}(\lambda a_1)\)；于是

\[
\left[ a _ {1}, a _ {2} \right] = (\operatorname{ad} a _ {1}). a _ {2} \in \mathfrak {h}
\]

由于 \(\mathfrak{h}\) 在 \(\mathrm{L}(\mathrm{G})\) 中闭，所以我们证明了 (i)。在证明的其余部分，固定一个李子群芽 \(\mathrm{H}'\)，它属于 \(\mathrm{G}\) 且李代数为 \(\mathfrak{h}\)。

设 \(V, v_0, f\) 如 (iii) 中。设 \(Y\) 是 \(G\) 中与 \(\mathfrak{h}\) 相关的左叶状结构（第1号）。对所有 \(y \in H'\)，\(T_y(H') = y\mathfrak{h}\)。另一方面，对所有 \(v \in V\)，\(T_v(V)\) 在 \(T_v(f)\) 下的像包含于 \(\mathfrak{h}_{f(v)} = f(v)\mathfrak{h}\)（按 \(\mathfrak{h}_{f(v)}\) 的定义）。根据《可微与解析流形》R，9.3.2，\(f\) 是从 \(V\) 到 \(Y\) 的态射。由于 \(H'\) 是 \(Y\) 的一片叶（《可微与解析流形》R，9.2.8），\(f(v) \in H'\) 当 \(v\) 充分接近 \(v_0\) 时。

设 \((a_{1},\ldots ,a_{s})\) 是 \(\mathfrak{h}\) 的一组基。存在 K 中 0 的一个开邻域 I 及映射 \(f_{1},\ldots ,f_{s}\)，它们是从 I 到 G 的 \(\mathbf{C}^r\) 类映射，满足 \(f_{j}(0) = e\)、\(f_{j}(\mathrm{I})\subset \mathrm{H}\)、\((\mathrm{Tf}_j)1 = a_j\)，对所有 \(j\)。根据 (iii)，\(f_{j}(\lambda)\in \mathrm{H}'\) 当 \(|\lambda |\) 充分小时。因此，\(f_{1}(\lambda_{1})f_{2}(\lambda_{2})\ldots f_{s}(\lambda_{s})\) 在 \(|\lambda_1|,\left|\lambda_2\right|,\ldots ,|\lambda_s|\) 充分小时构成 \(e\) 在 \(\mathbf{H}'\) 中的一个邻域；且这个邻域包含于 \(\mathbf{H}\)。所以得到 (iv)。

若 \(a \in \mathfrak{h}\)，则存在 0 的一个开邻域 \(\mathbf{I}\)，它位于 \(\mathbf{K}\) 中；以及一个映射 \(f\)，它是 \(\mathbf{C}^{\omega}\) 类的，从 \(\mathbf{I}\) 到 \(\mathbf{G}\)，使得 \(f(0) = e, f(\mathbf{I}) \subset \mathbf{H}'\)，\((\mathrm{T}_0f)1 = a\)。结合 (iv)，这推出 (v)。

定义 2. 称 \(\mathfrak{h}\) 为 H 在 \(e\) 处的切子代数。

命题 9. 设 G 是有限维李群，H 是 G 的子群。

\begin{enumerate}
\def\labelenumi{(\roman{enumi})}
\item
  H 上存在唯一的解析流形结构，具有如下性质：对 1 与 \(\omega\) 之间的任意 r、任意 \(C^r\) 类流形 V 及任意从 V 到 H 的映射 f，\(f\) 作为从 V 到 H 的映射属于 \(C^r\) 类，当且仅当 f 作为从 V 到 G 的映射属于 \(C^r\) 类。
\item
  在此结构下，\(\mathbf{H}\) 是李群，其典范嵌入 \(i\) 把 \(\mathbf{H}\) 映入 \(\mathbf{G}\)，是浸入，并且 \(\mathbf{L}(i)(\mathbf{L}(\mathbf{H}))\) 是 \(e\) 处 \(\mathbf{H}\) 的切李子代数。
\end{enumerate}

在 (i) 中，唯一性显然。我们证明存在性。设 \(\mathfrak{b}\) 是 H 在 \(e\) 处的切李代数。设 H' 是 G 的李子群芽，其李代数为 \(\mathfrak{b}\)。以 H' 的一个开子群芽替换 H' 后，可假设 H' ⊆ H（引理4 (iv)）。对所有 \(x \in \mathrm{H}\)，\(xH'x^{-1}\) 是 G 的李子群芽，其李代数为 \(x \mathfrak{h}x^{-1} = \mathfrak{b}\)。因此 H' ∩ \((xH'x^{-1})\) 在 H' 中开（第2号定理3），且映射 y → \(xyx^{-1}\) 是从 H' ∩ \(x^{-1}H'x\) 到 \(xH'x^{-1}\) ∩ H' 的同构。利用 § 1 第9号命题18，存在 H' 的一个开李子群芽 W 以及 H 上的一个李群结构，具有如下性质：W 在 H 中开，并且 H 与 H' 上的流形结构在 W 上诱导出相同的结构。于是，H 到 G 的典范嵌入 i 是浸入，且 L(i)(L(H)) = L(H') = \(\mathfrak{b}\)。此外，设 V 和 f 如 (i) 中。若 f: V → H 属于 \(C^r\) 类，则 \(i \circ f\): V → G 属于 \(C^r\) 类。反之，假设 \(i \circ f\): V → G 属于 \(C^r\) 类；我们证明 f: V → H 属于 \(C^r\) 类。通过平移，只需考虑存在 v0 ∈ V 使得 f(v0) = e，并证明 f: V → H 在 v0 的某个开邻域上属于 \(C^r\) 类的情形。现在，根据引理4 (iii)，当 v 充分接近 v0 时，f(v) ∈ H'，所述断言随即成立。因此我们证明了 (i)，而 (ii) 也在此过程中得到。

定义 3. 命题9中定义在 H 上的李群结构称为由 G 上的李群结构在 H 上诱导的结构。

若 H 是 G 的李子群，则 H 的李群结构由 G 上的李群结构诱导（《可微与解析流形》，R，5.8.5）。

若 \(\mathbf{G} = \mathbf{R}\) 且 \(\mathbf{H} = \mathbf{Q}\)，则 \(\mathfrak{h} = \{0\}\)，因而 \(\mathbf{H}\) 上的诱导结构是离散李群结构。同样，若 \(\mathbf{G} = \mathbf{C}\)（看作复李群）且 \(\mathbf{H} = \mathbf{R}\)，结论也成立。
% semantic-fragments: legacy-input-seam
\relax{}
